\chapter{相关研究}

本章围绕 GNSS NLOS 信号识别这一核心任务，梳理了相关技术的发展脉络，对国内外代表性研究工作进行总结。首先介绍基于几何模型和传统机器学习的检测方法，然后聚焦于基于深度学习的方法，分析各类方法的优势与不足。

\section{基于几何约束的检测方法}

早期 NLOS 检测主要依赖几何约束和统计特性。这类方法无需训练数据，实现简单，但泛化能力有限。

\subsection{卫星高度角约束}

卫星高度角约束是最直接的 NLOS 检测方法之一。低高度角卫星信号传播路径更长，受大气延迟和多路径效应影响更显著，且更容易被建筑物遮挡。因此，设置高度角阈值（如 15°或 30°）剔除低仰角卫星观测。

该方法计算开销极低，可直接嵌入接收机固件。然而，固定阈值难以适应不同城市环境：在高密度建筑群中，即使是中等高度角的卫星也可能被遮挡；而在开阔区域，过低的高度角阈值会浪费可用观测。

\subsection{一致性检测}

一致性检测通过比较多个接收机或多个历元的观测值来判断信号传播路径。若某卫星的伪距观测值与其他卫星或历史观测显著不一致，则可能为 NLOS 信号。

这类方法物理意义明确，但依赖冗余观测或多接收机配置，在单接收机场景下适用性受限。

\section{基于传统机器学习的检测方法}

随着机器学习技术的发展，研究者开始采用数据驱动的方法进行 NLOS 检测。传统机器学习方法可融合多维度特征进行分类，性能优于简单阈值法。

\subsection{支持向量机方法}

支持向量机（Support Vector Machine, SVM）是早期应用广泛的分类方法。Hsu 等人（2015）提出基于 SVM 的 NLOS 检测方法，融合载噪比、高度角、方位角和伪距残差等特征，在城市场景下取得约 70% 的分类准确率。

SVM 方法的优势在于小样本下表现稳定、可解释性较好。然而，核函数选择和参数调优对性能影响显著，且难以处理大规模数据集。

\subsection{随机森林方法}

随机森林（Random Forest）通过集成多棵决策树提高分类性能。Wang 等人（2017）采用随机森林进行 NLOS 识别，通过特征重要性分析发现载噪比和高度角是最具判别性的特征。

相比 SVM，随机森林对参数选择更鲁棒，且可处理非线性特征交互。但特征工程仍依赖领域知识，特征选取策略对最终性能影响较大。

\subsection{梯度提升树方法}

梯度提升树（Gradient Boosting Decision Tree, GBDT）及其变体（如 XGBoost、LightGBM）在多个分类任务上取得优异性能。Zhang 等人（2020）采用 XGBoost 进行 NLOS 检测，在相同特征下较随机森林提升约 3 个百分点。

然而，传统机器学习方法普遍存在以下局限：首先，依赖人工特征工程，难以捕捉特征间的复杂时空耦合关系；其次，模型容量有限，难以充分学习大规模数据的判别性表示。

\section{基于深度学习的检测方法}

近年来，深度学习在计算机视觉、自然语言处理等领域取得突破性进展，研究者开始将深度神经网络引入 NLOS 检测任务。

\subsection{全连接神经网络}

全连接神经网络（Fully Connected Neural Network, FCNN）是最早应用于 NLOS 检测的深度学习模型。Chen 等人（2018）设计了一个三层全连接网络，输入为单历元的多维度特征，输出为 LOS/NLOS 分类概率。

FCNN 可自动学习特征的非线性组合，避免手工特征设计。然而，单历元输入忽略了卫星信号的时序演化信息，难以捕捉 NLOS 状态的时序连续性。

\subsection{循环神经网络方法}

循环神经网络（Recurrent Neural Network, RNN）及其变体 LSTM、GRU 可建模时序依赖关系。Li 等人（2020）采用双向 LSTM 处理时间窗口内的 GNSS 观测序列，在动态场景下准确率达到 75%。

RNN 类方法的优势在于可捕捉信号的时序演化模式，但计算复杂度较高，且双向 LSTM 利用未来信息进行预测，不符合实时检测场景的因果约束。

\subsection{注意力机制方法}

注意力机制（Attention Mechanism）允许模型动态关注输入序列中的重要部分。Transformer 架构在序列建模任务上取得突破性进展。Wu 等人（2022）将 Transformer 引入 NLOS 检测，通过自注意力机制建模卫星间的空间关系。

然而，现有深度学习方法普遍存在以下不足：首先，多数方法仅利用单历元空间特征或单卫星时序特征，未能同时建模空间分布和时间演化；其次，跨场景泛化能力不足，在一个地点训练的模型直接部署到另一地点时性能显著下降；此外，模型复杂度与实时性之间的平衡尚未得到充分解决。

图~\ref{fig:method-taxonomy} 展示了 GNSS NLOS 检测方法的分类体系。

\begin{figure}[htbp]
\centering
\includegraphics[width=0.9\textwidth]{figures/gnss_nlos_method_taxonomy.png}
\caption{GNSS NLOS 检测方法分类体系}
\label{fig:method-taxonomy}
\end{figure}

图~\ref{fig:research-framework} 展示了本文 STCA 方法的研究框架。

\begin{figure}[htbp]
\centering
\includegraphics[width=0.9\textwidth]{figures/stca_research_framework.png}
\caption{本文 STCA 方法研究框架图}
\label{fig:research-framework}
\end{figure}

\section{本章小结}

本章系统综述了 GNSS NLOS 信号检测领域的国内外研究进展，包括基于几何约束的检测方法、基于传统机器学习的方法和基于深度学习的方法。各类方法在不同程度上推动了该领域的发展，但仍存在未能同时利用时空特征、跨场景泛化能力不足、实时性受限等问题。针对这些不足，下一章将详细介绍本文提出的 STCA 模型设计。
